% ---------------------------------------------------------------------------
% Fachmodul Reguläre Sprachen (Kristian Popov)
%
% Inhaltlich abgeleitet aus Abschnitt 4.3 der Dokumentation.
%
% Der ausformulierte Sprechtext liegt in:
%     docs/presentation/Sprechtext_Automaten_Sprachen.md
%
% Die Kurznotizen stehen in \note{...}. Sie sind standardmäßig unsichtbar.
% Zum Anzeigen in presentation.tex vor \begin{document} ergänzen:
%     \setbeameroption{show notes}
% oder für einen zweiten Bildschirm:
%     \setbeameroption{show notes on second screen=right}
% ---------------------------------------------------------------------------

\addsectionframe{Fachmodul: Reguläre Sprachen}

% ---------------------------------------------------------------------------
\begin{fullframe}{Reguläre Sprachen -- drei Arbeitsweisen}
  \begin{center}
    \begin{tikzpicture}[
        >=Stealth,
        box/.style={draw=thwsblue, very thick, rounded corners=3pt, fill=thwsblue!5,
                    minimum width=30mm, minimum height=12mm, align=center, font=\small},
        res/.style={draw=thwspetrol, very thick, rounded corners=3pt, fill=thwspetrol!8,
                    minimum width=26mm, minimum height=10mm, align=center, font=\small},
        lbl/.style={font=\scriptsize, thwsorange}
      ]
      \node[box] (gram) {Grammatik\\{\scriptsize rechtslinear}};
      \node[box, right=42mm of gram] (lang) {Sprach-\\beschreibung};

      \draw[->, very thick, thwsorange] ([yshift=3mm]gram.east)
        -- node[lbl, above, align=center] {Aufzählung\\{\scriptsize begrenzt: Länge, Anzahl}} ([yshift=3mm]lang.west);
      \draw[<-, very thick, thwsorange] ([yshift=-5mm]gram.east)
        -- node[lbl, below, align=center] {Erzeugung\\{\scriptsize strikt rechtslinear}} ([yshift=-5mm]lang.west);

      \node[res, below=16mm of gram] (word) {Wort + Grammatik};
      \node[res, right=42mm of word] (verdict) {akzeptiert / abgelehnt\\{\scriptsize mit Begründung}};
      \draw[->, very thick, thwspetrol] (word) -- node[lbl, above, thwspetrol] {Worterkennung} (verdict);
    \end{tikzpicture}
  \end{center}

  \vspace{-1mm}
  \footnotesize
  Alle drei Modi bleiben im Hintergrund erhalten und werden nur ein- oder ausgeblendet --
  dadurch überlebt eine begonnene Eingabe den Moduswechsel.
\end{fullframe}
\note{%
  Kurz durch die drei Modi gehen, nicht vertiefen.
  Der Satz unten ist eine UI-Entscheidung: Komponenten bleiben gemountet, damit Eingaben nicht verloren gehen.
  Überleitung zur nächsten Folie: "Nicht jede Eingabe wird akzeptiert -- und das mit Absicht."
}

% ---------------------------------------------------------------------------
\begin{fullframe}{Bewusst abgelehnte Sprachbeschreibungen}
  \begin{center}
    \begin{tikzpicture}[
        bad/.style={draw=red!65!black, very thick, rounded corners=3pt, fill=red!5,
                    minimum width=52mm, minimum height=13mm, font=\small\ttfamily},
        why/.style={font=\footnotesize, align=left, text width=58mm}
      ]
      \node[bad]                     (a) {L = \{a\^{}n b\^{}n | n > 0\}};
      \node[why, right=6mm of a]     (aw) {\textbf{nicht regulär} -- dieselbe Variable in zwei
                                            Faktoren koppelt die Wortlängen aneinander};

      \node[bad, below=9mm of a]     (b) {L = \{a\^{}n | n < 5\}};
      \node[why, right=6mm of b]     (bw) {\textbf{obere Schranke} -- beschreibt zwar eine endliche
                                            Sprache, ist als Aufzählung aber klarer};
    \end{tikzpicture}
  \end{center}

  \vspace{1mm}
  \footnotesize
  Beide Eingaben ließen sich formal verarbeiten. Sie werden trotzdem mit einer
  \textbf{eigenen Meldung} zurückgewiesen: Diese Fälle ausdrücklich zu erklären ist lehrreicher,
  als sie stillschweigend zu verwerfen.
\end{fullframe}
\note{%
  Hier ist die didaktische Begründung der Punkt, nicht die Technik.
  a\^{}n b\^{}n ist das Standardbeispiel für eine nicht-reguläre Sprache -- gut geeignet, um zu zeigen,
  dass das Werkzeug die Theorie ernst nimmt.
  Mögliche Rückfrage: "Warum keine oberen Schranken?" -- Antwort: endlich, also regulär, aber die
  Aufzählung ist die klarere Schreibweise; die Eingabesprache bleibt kleiner.
}

% ---------------------------------------------------------------------------
\begin{fullframe}{Übungsprüfung: Verhalten statt Struktur}
  \begin{center}
    \begin{tikzpicture}[
        >=Stealth,
        aut/.style={draw=thwsblue, very thick, rounded corners=3pt, fill=thwsblue!5,
                    minimum width=34mm, minimum height=14mm, align=center, font=\small},
        test/.style={draw=thwsorange, very thick, rounded corners=3pt, fill=thwsorange!8,
                     minimum width=40mm, minimum height=14mm, align=center, font=\small},
        ok/.style={font=\footnotesize, thwspetrol}
      ]
      \node[aut]                    (stud) {Lösung der\\Studierenden};
      \node[aut, below=8mm of stud] (ref)  {Referenzlösung};
      \node[test, right=26mm of $(stud)!0.5!(ref)$] (words)
        {alle Wörter über dem\\gemeinsamen Alphabet\\{\scriptsize bis zur festgelegten Länge}};

      \draw[->, very thick, thwsblue] (stud.east) -- (words.west);
      \draw[->, very thick, thwsblue] (ref.east)  -- (words.west);

      \node[align=left, font=\footnotesize, right=6mm of words, text width=46mm]
        {gleiche Entscheidung für jedes Wort? \\[1mm]
         \textcolor{thwspetrol}{\textbf{ja}} $\Rightarrow$ bestanden \\
         \textcolor{red!65!black}{\textbf{nein}} $\Rightarrow$ erstes abweichendes Wort als Gegenbeispiel};
    \end{tikzpicture}
  \end{center}

  \vspace{1mm}
  \footnotesize
  Ein \textbf{Strukturvergleich} würde korrekte Lösungen ablehnen, die nur anders aufgebaut sind.
  Die begrenzte Wortmenge ist bewusst eine \textbf{didaktische Rückmeldung} und kein formaler
  Äquivalenznachweis. Strukturelle Mängel werden vorrangig gemeldet -- sie sind meist die eigentliche Ursache.
\end{fullframe}
\note{%
  Wichtig: das ist eine bewusste Grenze, kein Mangel. Offen sagen, dass es kein Äquivalenzbeweis ist.
  Der Standardwert für die Wortlänge ist vier; einzelne Aufgaben setzen ihn höher.
  Das Gegenbeispiel ist der didaktische Mehrwert gegenüber einem blossen "falsch".
}

% ---------------------------------------------------------------------------
\begin{fullframe}{Wiederverwendung statt Duplikat}
  \begin{center}
    \begin{tikzpicture}[
        >=Stealth,
        mod/.style={draw=thwsblue, very thick, rounded corners=3pt, fill=thwsblue!5,
                    minimum width=44mm, minimum height=9mm, font=\small\ttfamily},
        rl/.style={draw=thwspetrol, very thick, rounded corners=3pt, fill=thwspetrol!8,
                   minimum width=44mm, minimum height=9mm, font=\small\ttfamily},
        hd/.style={font=\footnotesize\bfseries}
      ]
      \node[hd] (h1) {Reguläre Sprachen};
      \node[rl, below=3mm of h1]  (r1) {rlLanguageOps};
      \node[rl, below=2mm of r1]  (r2) {rlGrammarValidation};

      \node[hd, right=34mm of h1] (h2) {Endliche Automaten};
      \node[mod, below=3mm of h2] (m1) {faMachine};
      \node[mod, below=2mm of m1] (m2) {faGrammarConversion};
      \node[mod, below=2mm of m2] (m3) {faGrammarParsing};

      \draw[->, very thick, thwsorange] (r1.east) -- (m1.west);
      \draw[->, very thick, thwsorange] (r1.east) -- (m2.west);
      \draw[->, very thick, thwsorange] (r2.east) -- (m3.west);

      \node[font=\scriptsize\itshape, gray, below=4mm of m3, align=center]
        {Abhängigkeit nur in eine Richtung -- kein Zyklus};
    \end{tikzpicture}
  \end{center}

  \vspace{-3mm}
  \footnotesize
  Das Verzeichnis \texttt{domain} hat \textbf{keine Abhängigkeit zu React} und bleibt ohne
  gerenderte Oberfläche testbar. Die Aufzählung führt eine Grammatik über denselben
  Automatenkern aus wie die Simulation -- eine eigene Ableitungslogik hätte zwei
  Implementierungen desselben Verhaltens erfordert.
\end{fullframe}
\note{%
  Abschluss des Blocks: die Trennung Fachlogik/Oberfläche hat sich konkret ausgezahlt.
  Guter Anknüpfungspunkt für Fragen zu Tests -- die Domain-Logik ist ohne UI testbar,
  darauf lässt sich der Testteil der Präsentation aufbauen.
}
